\section{An unramified nonabelian extension of the source field}
\label{sec:unramified}

\subsection{The degree-three extension costs no new different}

Let $k=\Q(\sqrt{229})$. It is the quadratic subfield of the $S_3$ splitting field $F$. The elementary ramification calculation above gives a particularly favorable arithmetic situation.

\begin{theorem}[Unramified cubic base change]\label{thm:unramified}
The extension $F/k$ is cyclic of degree three and unramified at every finite prime. Let $q$ be divisible by $229$ and let $K$ be the imaginary quadratic field from the source. Set
\[
 E_0=K\Q(\zeta_q),\qquad E=F E_0.
\]
Then $[E:E_0]=3$, this extension is unramified at every finite prime, and
\begin{equation}\label{eq:different-inherited}
 \DD_{E/K}=\DD_{E_0/K}\OO_E,\qquad
 \DD_{E/\Q}=\DD_{E_0/\Q}\OO_E,\qquad
 \rd(E)=\rd(E_0).
\end{equation}
\end{theorem}

\begin{proof}
The subgroup fixing $k$ is $A_3$, so $F/k$ is cyclic cubic. Away from $229$ the splitting field is unramified. At $229$ its inertia group is a transposition subgroup of $S_3$, whose intersection with $A_3$ is trivial. Relative inertia for $F/k$ is exactly that intersection. Thus the relative different is the unit ideal.

The quadratic field $k$ has conductor $229$, so it lies in $\Q(\zeta_q)$. The field $E_0/\Q$ is abelian. Therefore $F\cap E_0$ is an abelian Galois subextension of $F/\Q$, and is contained in $k$; since it contains $k$, equality holds. Hence $[E:E_0]=3$.

An unramified extension of local fields remains an unramified finite etale algebra after base change; each field factor is unramified. Apply this at all finite places to $F/k$. This gives $\DD_{E/E_0}=\OO_E$. Transitivity of the different gives the first two identities in \eqref{eq:different-inherited}. The relative discriminant formula gives
\[
 |\operatorname{disc}(E)|=|\operatorname{disc}(E_0)|^3,
\]
and $[E:\Q]=3[E_0:\Q]$, proving the root-discriminant identity.
\end{proof}

\begin{figure}[htbp]
\centering
\begin{tikzpicture}[>=Stealth,every node/.style={align=center}]
\node (E) at (0,2.8) {$E=F E_0$};
\node (E0) at (-2.8,1.35) {$E_0=K\Q(\zeta_q)$};
\node (F) at (2.8,1.35) {$F$};
\node (k) at (0,0) {$k=\Q(\sqrt{229})$};
\draw (E0) -- node[above=3pt,midway,font=\small,sloped] {degree $3$} (E);
\draw (F) -- (E);
\draw (k) -- (E0);
\draw (k) -- node[above=3pt,midway,font=\small,sloped] {degree $3$} (F);
\end{tikzpicture}
\caption{The favorable nonabelian benchmark, with $F\cap E_0=k$. Both displayed degree-three extensions are unramified at every finite prime.}
\end{figure}

This corrects the general different-related uncertainty in the earlier synthesis for its specific example. A principal generator of the old different remains a generator after extension. This statement does not say that an arbitrary fractional ideal is principal, or that adjoining an arbitrary nonabelian field has the same property.

\subsection{The fractional lattice and trace self-duality}

From this point on, let $q$ be an admissible source tower modulus, with $229$ included in its fixed base modulus. Let $A_0$ be the source's fractional Weil lattice. In ideal notation its local construction gives
\begin{equation}\label{eq:half-different}
 A_0\overline{A_0}=\DD_{E_0/K}^{-1}.
\end{equation}
It permits inverse-different denominators on one side of each selected conjugate pair of primes, and remains integral on the other side. This is described explicitly by
\src{New/Derivatives/Vacuum/ShortVectors}{finiteImage_mem_iff} and the finite-idele valuation identity in
\src{New/Derivatives/Field/WeilData}{v_alpha_mul_star_mul_relDifferentGen}.\cite{shortVectors,weilField}

For the positive definite Hermitian pairing
\[
 h(x,y)=\Tr_{E/K}(\overline{x}y),
\]
the dual of a fractional ideal $A$ is
\begin{equation}\label{eq:hermitian-dual}
 A^\#=\DD_{E/K}^{-1}\overline A^{-1}.
\end{equation}
Indeed, $h(z,A)\subseteq\OO_K$ is equivalent to
$\overline z A\subseteq\DD_{E/K}^{-1}$; conjugate this condition and use conjugation stability of the different. Trace duality and different transitivity are standard number-field facts.\cite{Milne}

\begin{corollary}[The corrected lattice]
The ideal $A=A_0\OO_E$ satisfies
\[
 A\overline A=\DD_{E/K}^{-1},\qquad A^\#=A.
\]
It is a full projective $\OO_K$-lattice in $E$, with its positive definite Hermitian trace pairing.
\end{corollary}

\begin{proof}
Extend \eqref{eq:half-different} to $\OO_E$ and apply Theorem~\ref{thm:unramified}. Equation~\eqref{eq:hermitian-dual} gives self-duality. A fractional ideal is finite and torsion-free over the Dedekind domain $\OO_K$, hence projective. Under the induced CM type, the trace of $\overline x x$ is a sum of squared complex absolute values, proving positivity.
\end{proof}

For clarity, a more general ideal construction is also available. If every rational prime below $\DD_{E/K}$ splits in $K$, choose one of its two $K$-primes. Put valuation $-v_{\mathfrak P}(\DD_{E/K})$ in $A$ above the chosen prime and valuation zero above its conjugate. Then $A\overline A=\DD_{E/K}^{-1}$ prime by prime. If $E/K$ is Galois, this choice is stable under that Galois group. No principal different is required. The special unramified benchmark is stronger because it reuses the old lattice directly.

\subsection{An existing source theorem supplies the Weil data}

An important internal theorem was absent from the earlier transfer discussion:
\begin{center}
\src{Lib/ClassField/Grunwald}{New.nonempty_weilData_of_not_dvd_differentIdeal}.
\end{center}
It says that a CM extension with compatible CM types inherits Weil data when no finite prime fixed by conjugation divides the relative different. Its proof retains the purely imaginary parameter $j$, pulls the splitting character along the idele norm, and constructs a compatible finite idele.\cite{grunwald}

Theorem~\ref{thm:unramified} verifies the hypothesis in the strongest possible way. Here one can retain the old finite idele by extension, because the different is unchanged. The remaining structural conditions also transport: the root of unity of order $5040$ already lies in the base, the parameter $j$ has the same positive imaginary values on the induced CM type, and its trace-one algebraic integer remains integral. The splitting character is the original one composed with the idele norm $E\to E_0$.

This is an application of an existing theorem in the proof library. It must not be counted as a newly proved general theorem or as a newly compiled Lean instantiation.

\subsection{Period normalization and the old unit witnesses}

The source frame construction needs more than the matrix identity. For each $p$-adic embedding, at least one frame member $y$ must make all of its normalized CM-coordinate coefficients units. In the source, the two possible witnesses are
\begin{equation}\label{eq:unit-witnesses}
 1,\qquad
 \zeta_q\,\mathrm{dualVector}(\varphi(q)-1).
\end{equation}
Their property is the norm assertion
\[
 \left|\iota_p\bigl(c\,\overline{\tau(y)}\bigr)\right|_p=1
 \quad\text{for every CM embedding }\tau,
\]
with the source's period normalization $c$.\cite{shortVectors,tightFrames}

Every new CM embedding restricts to an old one. The finite-idele valuations and root discriminant extend unchanged. Consequently the same norm assertions hold for the two old witnesses in $A$.

There is a normalization subtlety. The source defines its vacuum root by a principal complex root. Under degree-three base change, the quantity whose root is taken is cubed and the root degree triples. Recomputing the principal root need not give the same complex number: the new and old choices can differ by a root of unity. Their absolute values at all archimedean and $p$-adic places agree. Alternatively, choose the inherited old root as the new normalization; it satisfies the new power equation. Thus the unit-witness statement transports as an equality of norms. Literal equality of independently chosen principal-root values is not asserted.

The witnesses alone cannot span the new embedding directions. The next section proves that one can include them in genuine frames at every sufficiently large integer bound, without assuming a free lattice.
